\section*{第 \ref{ch:g11:enzymes-and-phenotype} 章——酶与表型}

\begin{solution}{exo:g11:enzymes-and-phenotype:1}
\omterm{def:g11:enzymes-and-phenotype:phenotype}{基因型}：携带的\omterm{def:g10:universal-dna:gene}{等位基因}的集合。\omterm{def:g11:enzymes-and-phenotype:phenotype}{表型}：可观测的
特征的集合，分别在分子水平、\omterm{def:g10:cells-common-unit:cell}{细胞}水平和个体水平上。
\end{solution}

\begin{solution}{exo:g11:enzymes-and-phenotype:2}
酶上\omterm{def:g11:enzymes-and-phenotype:enzyme}{底物}结合并被转化的口袋，即\omterm{def:g11:enzymes-and-phenotype:enzyme}{活性位点}。
专一性：该位点只与一种\omterm{def:g11:enzymes-and-phenotype:enzyme}{底物}（或一类\omterm{def:g11:enzymes-and-phenotype:enzyme}{底物}）契合，
并催化一种反应。
\end{solution}

\begin{solution}{exo:g11:enzymes-and-phenotype:3}
接近 \qty{37}{\celsius}；在 \qty{45}{\celsius} 时约剩三分之一。
\end{solution}

\begin{solution}{exo:g11:enzymes-and-phenotype:4}
分子水平：酪氨酸酶缺失或失活，不能合成黑色素。\omterm{def:g10:cells-common-unit:cell}{细胞}水平：
色素\omterm{def:g10:cells-common-unit:cell}{细胞}存在，但其中没有黑色素。个体水平：毛发白色，
皮肤和眼睛极浅，对光敏感，易被晒伤。
\end{solution}

\begin{solution}{exo:g11:enzymes-and-phenotype:5}
胃蛋白酶的折叠构象仅在强酸（pH 2）中起作用；肠道
呈中性至弱碱性，曲线显示此时活性为零。
\end{solution}

\begin{solution}{exo:g11:enzymes-and-phenotype:6}
一个有功能的\omterm{def:g10:universal-dna:gene}{等位基因}产生正常酶量的一半，但
每个酶分子每秒可工作数千次：即使只有一半的量，
仍足以转化出足够的酪氨酸以合成正常色素。该\omterm{def:g11:enzymes-and-phenotype:phenotype}{表型}
已达饱和；只有两个失效的\omterm{def:g10:universal-dna:gene}{等位基因}才会表现出来。
\end{solution}

\begin{solution}{exo:g11:enzymes-and-phenotype:7}
双亲的白化病源于两个不同\omterm{def:g10:universal-dna:gene}{基因}的阻断（两种不同的
酶，分属\omterm{prop:g11:enzymes-and-phenotype:pathway}{代谢途径}或黑色素转运过程）。孩子从阻断位置
不同的那一方亲本各继承一个有功能的\omterm{def:g10:universal-dna:gene}{等位基因}，
因此两种酶都能工作，色素得以合成。
\end{solution}

\begin{solution}{exo:g11:enzymes-and-phenotype:8}
苯丙氨酸积累；酪氨酸（及由其衍生的物质）
缺乏。黑色素由酪氨酸合成，因此其短缺会减少
色素：未经治疗的苯丙酮尿症患者中，浅色毛发和皮肤较为常见。
\end{solution}

\begin{solution}{exo:g11:enzymes-and-phenotype:9}
在以乳汁为主食的地区，消化乳汁可提供能量和钙，
且不致患病，因此持续型\omterm{def:g10:universal-dna:gene}{等位基因}具有优势。在
成人不饮用乳汁的地方，该\omterm{def:g10:universal-dna:gene}{等位基因}并无作用；它
只有在提供\omterm{def:g11:enzymes-and-phenotype:enzyme}{底物}的环境中才算“更健康”。
\end{solution}

\begin{solution}{exo:g11:enzymes-and-phenotype:10}
就酪氨酸酶而言：分子水平上\omterm{def:g11:enzymes-and-phenotype:enzyme}{底物}不再契合，不能合成黑色素；
\omterm{def:g10:cells-common-unit:cell}{细胞}水平上色素\omterm{def:g10:cells-common-unit:cell}{细胞}内无黑色素；个体水平上为白化病。就乳糖酶而言：
乳糖未被消化；肠道\omterm{def:g10:cells-common-unit:cell}{细胞}缺少该糖，且结肠
对其\omterm{def:g10:cell-metabolism:fermentation}{发酵}；即乳糖不耐受。
\end{solution}

\begin{solution}{exo:g11:enzymes-and-phenotype:11}
其四肢温度高于 \qty{35}{\celsius}，酶几乎
不起作用，小猫全身都很白。冬季户外时，耳朵
和爪子变凉，酶变得活跃，在这些部位再长出的
皮毛呈深色。
\end{solution}

\begin{solution}{exo:g11:enzymes-and-phenotype:12}
分子水平：一半血红蛋白为镰刀形。\omterm{def:g10:cells-common-unit:cell}{细胞}水平：在
正常氧气供应下\omterm{def:g10:cells-common-unit:cell}{细胞}保持圆形；极低氧时，部分\omterm{def:g10:cells-common-unit:cell}{细胞}呈镰刀形。
个体水平：日常生活中健康，在高海拔地区或
潜水时有风险——但能抵抗疟疾。环境（氧气、
疟疾）决定携带者的\omterm{def:g11:enzymes-and-phenotype:phenotype}{表型}是负担还是
优势。
\end{solution}

\begin{solution}{exo:g11:enzymes-and-phenotype:13}
在 \qty{41}{\celsius} 时，酶仍保持大部分活性
（曲线接近峰值）；在 \qty{43}{\celsius} 时活性急剧下降，
此时\omterm{def:g11:gene-expression:protein}{蛋白质}开始展开；若展开过程继续，
则该过程不可逆：酶被破坏，而不只是活性减弱。
\end{solution}

\begin{solution}{exo:g11:enzymes-and-phenotype:14}
每条\omterm{def:g10:muscles-and-joints:muscle}{肌纤维}中\omterm{def:g10:cells-common-unit:organelle}{线粒体}和毛细血管更多；心脏更大，
每搏输出量更大；每升血液中血红蛋白更多；
\omterm{def:g10:exercise-and-energy:glycogen}{糖原}储备更多。这些都是同一\omterm{def:g11:enzymes-and-phenotype:phenotype}{基因型}
在不同环境下产生的\omterm{def:g11:enzymes-and-phenotype:phenotype}{表型}：“遗传”性状是
\omterm{def:g11:enzymes-and-phenotype:phenotype}{基因型}所允许的变化范围，而非固定值。
\end{solution}

\begin{solution}{exo:g11:enzymes-and-phenotype:15}
\omterm{def:g11:enzymes-and-phenotype:phenotype}{基因型}决定合成哪种酶及其行为：无酶
1（苯丙酮尿症）、无酪氨酸酶（白化病）、热敏型
酪氨酸酶（暹罗猫）。环境只能在酶允许的范围内改变结果：
例如饮食可去除\omterm{def:g11:enzymes-and-phenotype:enzyme}{底物}，从而保护
大脑，但无法产生缺失的酶；低温使暹罗猫的
酶发挥作用，但任何温度都无法使白化个体的酶发挥作用；
环境中没有任何因素能让白化个体获得色素。\omterm{def:g11:enzymes-and-phenotype:phenotype}{基因型}设定范围；
环境则在范围内做出选择。
\end{solution}

\begin{solution}{pb:g11:enzymes-and-phenotype:1}
\textbf{1.} 躯干 \qty{38}{\celsius}：约 2\%，浅色。腿部 
\qty{36}{\celsius}：约 20\%，浅色。爪子 \qty{33}{\celsius}：80\%，
深色。耳朵和口鼻部 \qty{31}{\celsius}：90\% 以上，深色。

\textbf{2.} 身体和腿部浅色，爪子、耳朵、口鼻部（及尾巴）深色：
即暹罗猫的花色模式。

\textbf{3.} 各处均为 100\%：全身色素均匀分布的猫。

\textbf{4.} 在子宫内，小猫各部位均处于
\qty{38}{\celsius}，酶始终无活性，出生前也未沉积任何色素
；出生后末端渐冷，
色点随之变深。

\textbf{5.} 约 \qty{35.5}{\celsius}，即活性降至
30\% 的临界处：33 至
\qty{37}{\celsius} 之间曲线陡降的一段界定了边界。

\textbf{6.} 深色。在 \qty{30}{\celsius} 时，酶
在生长中的毛发里完全有活性。

\textbf{7.} 浅色。在 \qty{38}{\celsius} 时，酶无活性。

\textbf{8.} 色素随毛发生长沉积于其中，之后既
不能移除，也不能再添加；颜色只有在毛发脱落、
由该处温度下新生的毛发取代时才会改变。

\textbf{9.} 斑块处\omterm{def:g10:cells-common-unit:cell}{细胞}的\omterm{def:g11:enzymes-and-phenotype:phenotype}{基因型}前后相同，
与躯干其他部位一致；改变的只是
酶所处的温度，连带色素也随之改变。作用的对象是
酶的活性，这是一个分子水平的事件。

\textbf{10.} 猫的所有\omterm{def:g10:cells-common-unit:cell}{细胞}均来自同一枚卵，携带相同的
\omterm{def:g10:universal-dna:gene}{等位基因}；剃毛斑块的实验只改变了
一个区域的颜色，而未改变其\omterm{def:g10:cells-common-unit:cell}{细胞}——\omterm{def:g10:universal-dna:gene}{等位基因}不会因
冰袋而切换。

\textbf{11.} 分子水平：酪氨酸酶仅在约
\qty{35}{\celsius} 以下有活性。\omterm{def:g10:cells-common-unit:cell}{细胞}水平：末端的色素\omterm{def:g10:cells-common-unit:cell}{细胞}充满黑色素，
躯干处则空无黑色素。个体水平：浅色的猫，带
深色的色点。

\textbf{12.} 该\omterm{def:g10:chemistry-of-life:families}{氨基酸}位于能稳定折叠结构的位置；
替换后的\omterm{def:g10:chemistry-of-life:families}{氨基酸}在低温下能维持折叠，
但温度稍高便会松散，因此\omterm{def:g11:enzymes-and-phenotype:enzyme}{活性位点}
比普通酶更早失去形状。

\textbf{13.} 普通\omterm{def:g10:universal-dna:gene}{等位基因}编码的酶在体温下完全有活性；
即使只有正常量的一半，仍绰绰有余，
足以在全身合成色素。暹罗\omterm{def:g10:universal-dna:gene}{等位基因}相对于它为隐性。

\textbf{14.} 分子水平：任何温度下均无有效的酶，
而暹罗猫的酶仅在特定条件下起作用。\omterm{def:g10:cells-common-unit:cell}{细胞}水平：全身均无黑色素，
而暹罗猫在低温区域有黑色素。个体水平：全白的猫，眼睛为
粉红色，而暹罗猫为色点花纹。低温对白化个体
毫无影响。

\textbf{15.} 它们携带的是同一\omterm{def:g10:universal-dna:gene}{基因}的
温度敏感型\omterm{def:g10:universal-dna:gene}{等位基因}；很可能是某种类似的替换
使酶变得不稳定，两个\omterm{def:g10:biodiversity-scales:species}{物种}各自独立产生。

\textbf{16.} 二十倍；酶 1，即将苯丙氨酸
转化为酪氨酸的酶。

\textbf{17.} 过量的苯丙氨酸专门干扰发育中的
脑\omterm{def:g10:cells-common-unit:cell}{细胞}（影响其他\omterm{def:g10:chemistry-of-life:families}{氨基酸}的供应及
其成熟）；皮肤\omterm{def:g10:cells-common-unit:cell}{细胞}则可耐受。分子缺陷处处相同，
而\omterm{def:g10:cells-common-unit:cell}{细胞}水平的后果却并非如此。

\textbf{18.} \num{1200} 的五分之一为 \qty{240}{\micro mol/L}，低于
\num{360}，是安全的。若无此酶，血液中的水平
便直接随摄入量变化，因此\omterm{def:g11:enzymes-and-phenotype:phenotype}{表型}由\omterm{def:g11:enzymes-and-phenotype:enzyme}{底物}供应决定——即
环境——而非仅由缺失的酶决定。

\textbf{19.} 苯丙氨酸是构建身体\omterm{def:g11:gene-expression:protein}{蛋白质}所必需的，
且不能自行合成；若完全缺乏，生长便会停止。摄入量
须恰好够\omterm{def:g11:gene-expression:protein}{蛋白质}合成之用，不能过量，这就需要测量
并加以调节。

\textbf{20.} \omterm{def:g11:enzymes-and-phenotype:phenotype}{基因型}决定了酶（在猫中为热敏感型，
在患儿中则缺失）；环境决定了其条件（皮肤
温度、饮食中的苯丙氨酸）；\omterm{def:g11:enzymes-and-phenotype:phenotype}{表型}则是
\omterm{def:g11:enzymes-and-phenotype:phenotype}{基因型}在特定环境中表达的结果。
\end{solution}
